\chapter{Manuel --- Directions et Responsables P\'edagogiques}
\label{chap:directions}

\begin{rolebox}{directeur-departement, vice-president-departement, directeur-adjoint-licence, responsable-service-pedagogique, responsable-adjoint-service-pedagogique, directeur-mention, directeur-specialite, responsable-formation, responsable-annee, directeur-etudes, responsable-qualite, responsable-international, referent-commun, directeur-adjoint-ecole}
  Ce chapitre couvre l'ensemble des r\^oles de direction au niveau d\'epartement, mention, parcours, niveau ou service. Ces r\^oles ont des droits d'\'ecriture limit\'es \`a leur p\'erim\`etre et ne peuvent pas g\'erer les utilisateurs.
\end{rolebox}

\section{Vue d'ensemble du profil}

Les r\^oles de direction et de responsabilit\'e p\'edagogique disposent des fonctionnalit\'es suivantes\,:

\begin{itemize}
  \item \textbf{Consultation de l'annuaire} (responsables, formations, structures, secr\'etariats).
  \item \textbf{Gestion des affectations} dans leur p\'erim\`etre (cr\'eation et modification de contacts de r\^ole).
  \item \textbf{G\'en\'eration d'organigrammes} pour leur structure et ses descendants.
  \item \textbf{Consultation de la biblioth\`eque} d'organigrammes.
  \item \textbf{Gestion des d\'el\'egations}.
  \item \textbf{Demandes de r\^oles} personnalis\'es.
  \item \textbf{Signalements}.
\end{itemize}

\begin{notebox}[Ce qui n'est PAS disponible]
  Ces r\^oles \textbf{ne peuvent pas}\,:
  \begin{itemize}
    \item Cr\'eer, modifier ou supprimer des comptes utilisateurs.
    \item Importer ou exporter des donn\'ees (sauf si la d\'el\'egation le pr\'evoit).
    \item G\'erer les ann\'ees acad\'emiques.
    \item Modifier les structures.
    \item Consulter le journal d'audit.
  \end{itemize}
\end{notebox}

\section{Tableau de bord}

\screenshot{img_dashboard_direction}{Tableau de bord -- Directions et Responsables}{Capturer le tableau de bord avec un compte de type direction (ex. directeur-departement ou directeur-mention). Montrer les tuiles disponibles pour ce profil, l'ann\'ee courante (non modifiable) et les notifications.}

\section{Annuaire -- Recherche}

L'annuaire est la fonctionnalit\'e principale pour les r\^oles de direction. Il permet de retrouver rapidement les responsables, formations, structures et secr\'etariats.

\begin{steps}
  \item Naviguer vers \textbf{Rechercher} dans le menu.
  \item S\'electionner l'onglet souhait\'e\,: \textbf{Responsables}, \textbf{Formations}, \textbf{Structures} ou \textbf{Secr\'etariats}.
  \item Saisir un texte dans la barre de recherche et/ou utiliser les filtres.
  \item Les r\'esultats s'affichent en dessous.
\end{steps}

\screenshot{img_annuaire_direction}{Annuaire -- Vue Directions}{Capturer la page Recherche avec l'onglet Responsables actif, une recherche effectu\'ee et les r\'esultats visibles. Montrer les filtres (composante, d\'epartement, mention) et le tableau de r\'esultats.}

Pour les d\'etails complets de l'annuaire, voir le chapitre~\ref{chap:commun}.

\section{Gestion des affectations dans le p\'erim\`etre}

Les r\^oles de direction peuvent g\'erer les affectations au sein de leur structure.

\subsection{Cr\'eer ou modifier une affectation}

\begin{steps}
  \item Naviguer vers \textbf{Responsables}.
  \item Rechercher la personne concern\'ee.
  \item Cliquer sur \textbf{Ajouter un r\^ole} (dans la fiche de la personne).
  \item S\'electionner le \textit{R\^ole}, la \textit{Structure} (limit\'ee au p\'erim\`etre autoris\'e).
  \item Cliquer sur \textbf{Enregistrer}.
\end{steps}

\begin{notebox}[P\'erim\`etre des affectations]
  Un directeur de d\'epartement ne peut cr\'eer des affectations que pour son d\'epartement et ses sous-structures (mentions, parcours, niveaux). Il ne peut pas affecter des r\^oles dans d'autres d\'epartements.
\end{notebox}

\screenshot{img_affectation_direction}{Cr\'eation d'affectation -- Direction}{Capturer le formulaire d'affectation pour un r\^ole de direction. Montrer que les menus d\'eroulants de structure sont limit\'es au p\'erim\`etre du r\^ole (ex. d\'epartements disponibles = uniquement le d\'epartement de l'utilisateur).}

\subsection{Mettre \`a jour un contact de r\^ole}

\begin{steps}
  \item Localiser l'affectation dans la liste \textbf{Responsables}.
  \item Cliquer sur \textbf{Contact de r\^ole}.
  \item Renseigner ou mettre \`a jour\,: \textit{Courriel fonctionnel}, \textit{T\'el\'ephone}, \textit{Bureau}.
  \item Cliquer sur \textbf{Enregistrer}.
\end{steps}

\begin{tipbox}
  La mise \`a jour des contacts de r\^ole est particuli\`erement importante pour les secr\'etariats et les responsables de formation, car ces informations sont visibles dans l'annuaire par tous les utilisateurs.
\end{tipbox}

\section{Fiches Structures}

Les directions peuvent consulter les fiches de structure dans leur p\'erim\`etre.

\begin{steps}
  \item Naviguer vers \textbf{Structures} ou rechercher dans l'annuaire.
  \item Cliquer sur la structure souhait\'ee.
  \item La fiche d\'etaill\'ee affiche\,: hi\'erarchie, responsables, secr\'etariat, sous-structures, statistiques.
\end{steps}

\screenshot{img_fiche_structure_direction}{Fiche structure -- vue Direction}{Capturer la fiche d\'etaill\'ee d'une structure (d\'epartement ou mention). Montrer les sections\,: titre, hi\'erarchie parente, liste des responsables avec leurs r\^oles, liste du secr\'etariat, sous-structures, et statistiques chiffr\'ees.}

\section{Organigrammes}

\subsection{G\'en\'erer un organigramme pour son p\'erim\`etre}

\begin{steps}
  \item Naviguer vers \textbf{Organigramme}.
  \item Choisir la vue (\textbf{Structures} ou \textbf{Personnes}).
  \item S\'electionner la structure racine parmi celles de son p\'erim\`etre.
  \item Appliquer les filtres (r\^ole, niveau, etc.).
  \item Cliquer sur \textbf{G\'en\'erer}.
\end{steps}

\screenshot{img_organigramme_direction}{Organigramme -- G\'en\'eration pour une Direction}{Capturer la page organigramme avec la vue Personnes, le s\'electeur de structure racine (montrant le d\'epartement ou la mention de l'utilisateur) et les filtres. Montrer le r\'esultat g\'en\'er\'e si possible.}

\subsection{Consulter la biblioth\`eque}

\begin{steps}
  \item Dans la page \textbf{Organigramme}, acc\'eder \`a la section \og Biblioth\`eque \fg{}.
  \item Filtrer par structure ou type de vue.
  \item Cliquer sur \textbf{Voir en structures} ou \textbf{Voir en personnes} pour consulter un organigramme.
\end{steps}

\begin{notebox}[Consultation \'etendue]
  Les r\^oles de direction peuvent consulter les organigrammes d\'ej\`a g\'en\'er\'es par d'autres utilisateurs, m\^eme pour des structures hors de leur p\'erim\`etre. En revanche, ils ne peuvent g\'en\'erer des organigrammes que pour leur p\'erim\`etre.
\end{notebox}

\section{D\'el\'egations}

Les directions peuvent d\'el\'eguer leurs droits \`a un autre utilisateur dans leur p\'erim\`etre.

\subsection{Cr\'eer une d\'el\'egation}

\begin{steps}
  \item Naviguer vers \textbf{D\'el\'egations}.
  \item Cliquer sur \textbf{Cr\'eer une d\'el\'egation}.
  \item S\'electionner le \textit{D\'el\'egu\'e} (utilisateur existant).
  \item S\'electionner la \textit{Structure concern\'ee} (dans votre p\'erim\`etre).
  \item Indiquer les \textit{dates de d\'ebut et de fin}.
  \item Cliquer sur \textbf{Enregistrer}.
\end{steps}

\screenshot{img_delegations_direction}{D\'el\'egations -- Directions}{Capturer la page D\'el\'egations avec la liste des d\'el\'egations actives et le formulaire de cr\'eation visible (ou le bouton Cr\'eer).}

\subsection{R\'evoquer une d\'el\'egation}

\begin{steps}
  \item Dans la liste, cliquer sur la d\'el\'egation.
  \item Cliquer sur \textbf{R\'evoquer}.
  \item Confirmer.
\end{steps}

\section{Demandes de r\^oles}

Les directions peuvent soumettre des demandes de r\^oles personnalis\'es pour des personnes de leur p\'erim\`etre.

\subsection{Soumettre une demande}

\begin{steps}
  \item Naviguer vers \textbf{Demandes de r\^oles}.
  \item Cliquer sur \textbf{Nouvelle demande}.
  \item Renseigner\,: \textit{Utilisateur}, \textit{R\^ole demand\'e}, \textit{Structure}, \textit{Justification}.
  \item Cliquer sur \textbf{Envoyer}.
\end{steps}

\screenshot{img_demandes_roles_direction}{Demandes de r\^oles -- Directions}{Capturer la page demandes de r\^oles avec la liste des demandes (en attente, approuv\'ees, refus\'ees) et le formulaire ou bouton de nouvelle demande.}

\subsection{Suivre ses demandes}

Les statuts possibles sont\,:
\begin{itemize}
  \item \textbf{En attente} : examen par les Services Centraux en cours.
  \item \textbf{Approuv\'ee} : r\^ole attribu\'e, affectation cr\'e\'ee.
  \item \textbf{Refus\'ee} : demande rejet\'ee (consulter le commentaire de justification).
\end{itemize}

\section{Signalements}

\begin{steps}
  \item Naviguer vers \textbf{Signalements}.
  \item Cliquer sur \textbf{Cr\'eer un signalement}.
  \item S\'electionner le \textit{type d'erreur} et la \textit{personne ou structure concern\'ee}.
  \item D\'ecrire le probl\`eme.
  \item Cliquer sur \textbf{Envoyer}.
\end{steps}

\screenshot{img_signalements_direction}{Signalements -- Directions}{Capturer la page Signalements. Montrer le formulaire de cr\'eation et la liste des signalements de l'utilisateur avec les statuts.}
